\documentclass[10pt,a4paper]{amsart}
\usepackage[margin=19mm]{geometry}
\usepackage{amsmath,amssymb,mathtools}
\usepackage[scr=rsfs,cal=euler]{mathalfa}
\usepackage{xfrac}
\usepackage[unicode,hidelinks,bookmarks=false]{hyperref}
\newcommand{\ie}{i.e.\ }
\newcommand{\cf}{cf.\ }
\newcommand{\resp}{respectively\ }
\newcommand{\Subsection}{\subsection}
\providecommand{\nobreakdash}{\nobreak\hbox{-}}
\setlength{\emergencystretch}{2em}
\multlinegap0pt
\usepackage{bm}
\usepackage{bbm}
\usepackage{dsfont}
\usepackage{xspace}
\usepackage{cancel}
\usepackage{mathtools}
\usepackage{amsthm}
\makeatletter
\@ifundefined{lemma}{\newtheorem{lemma}{Lemma}}{}
\makeatother
\newcommand\Pic{\mathrm{Pic\,}}
\newcommand\ra{\rightarrow}
\newcommand\wt{\widetilde}
\newcommand{\mb}[1]{\mathbb{#1}}
\newcommand{\mib}[1]{\textit{\textbf{#1}}}
\newcommand{\mr}[1]{\mathrm{#1}}
\title[Arithmetic Degrees are Cohomological Lyapunov Multipliers]{Arithmetic Degrees are Cohomological Lyapunov Multipliers}
\author{Jiarui Song, Junyi Xie, She Yang}
\date{Source: arXiv:2507.17643v1 (CC BY 4.0)}
\begin{document}
\maketitle

\noindent\emph{Source and licence.} Arithmetic Degrees are Cohomological Lyapunov Multipliers. Version arXiv:2507.17643v1 (CC BY 4.0), distributed under the Creative Commons Attribution 4.0 International licence, as recorded in the arXiv licence metadata for this exact version and shown on the abstract page https://arxiv.org/abs/2507.17643v1. Full rights notice: licenses/arxiv-2507.17643-CC-BY-4.0.txt.\par\smallskip

\noindent\textit{Editorial scope.} Counted unit: the Lemma whose source label is \texttt{Lem\_lift} in the article (source0.tex lines 357--359, source label \texttt{Lem\_lift}), delivered together with the article's own complete original proof of it (source0.tex lines 361--393, one \texttt{proof} environment of 33 lines). No mathematics is written, repaired or re-derived: statement and proof are the article's own text line for line. Required context retained uncounted: none. Every definition, display and label of those lines is retained unchanged. Omitted from this package and not redistributed: 1--356, 360--360, 394--535. Nothing inside a retained excerpt is omitted or altered: every retained body is delivered exactly as the article's lines are written, so no omission receipt of any kind accompanies this package. Citations: the retained excerpts carry no citation command at all, so no bibliography entry of the article is transcribed and no reference is redistributed. Reference adaptation: none was needed and none was made; every reference inside the delivered unit resolves inside it, so no unresolved reference or citation is produced. Printed slips kept literal: no printed word, symbol, hypothesis or number of the counted statement or of its proof was altered. Layout and class: the article's own class is \texttt{amsart} with packages including color, amsmath, hyperref, amscd, amsthm, amsfonts, amssymb, esint, fullpage, xy, mathtools, xcolor, mathrsfs, geometry; the wrapper is the lane's pinned \texttt{amsart} editorial wrapper at 10pt on A4, which restates the theorem-like environments the retained text needs and adds the article's own macro definitions that the retained text needs (\texttt{\textbackslash Pic}, \texttt{\textbackslash mb}, \texttt{\textbackslash mib}, \texttt{\textbackslash mr}, \texttt{\textbackslash ra}, \texttt{\textbackslash wt}), with extra wrapper packages (bm, bbm, dsfont, xspace, cancel, mathtools). The delivered numbering is the unit's own. Assets: no retained span contains a figure environment, a graphics-inclusion command, a PDF-page inclusion, a table or a TikZ picture; no image file is copied into this package, the package declares no asset dependency and no image file is redistributed with it. Licence: version arXiv:2507.17643v1 of this article is distributed under the Creative Commons Attribution 4.0 International licence, as recorded in the arXiv licence metadata for this exact version and shown on the abstract page https://arxiv.org/abs/2507.17643v1. A full rights notice is delivered at licenses/arxiv-2507.17643-CC-BY-4.0.txt. Characters: every retained excerpt is pure ASCII, so the delivered engine, XeLaTeX with its default Latin Modern text font, reports no missing character.
\begin{lemma}\label{Lem_lift}
Let $p: \mr{Pic}(X)_{\mb{R}}\ra \mr{N}^1(X)_{\mb{R}}$ be the natural projection. Let $\rho_0$ denote the spectral radius of the self-isogeny $f^*: \mr{Pic}^0(X)\ra \mr{Pic}^0(X)$. Then there exists a lifting $\iota: \bigoplus\limits_{\mu>\rho_0}E_{\mu}\ra \mr{Pic}(X)_{\mb{R}}$ such that $p\circ\iota$ is the natural inclusion and $\iota\circ f^*=f^*\circ\iota$. The last two $f^*$ are pull-backs on $\bigoplus\limits_{\mu>\rho_0}E_{\mu}\subseteq\mathrm{N}^1(X)_{\mathbb{R}}$ and $\mathrm{Pic}(X)_{\mathbb{R}}$, respectively.
\end{lemma}

\begin{proof}
It suffices to construct $\iota$ on each cyclic subspace of $E_\mu$ for $\mu > \rho_0$. Let $E \subset E_\mu$ be a cyclic subspace for $f^*: E_\mu \to E_\mu$, with basis $D_1, D_2, \ldots, D_r$ satisfying
\[f^*D_1\equiv \mu D_1+D_2, f^*D_2\equiv\mu D_2+D_3,\ldots, f^*D_r\equiv\mu D_r.\]

Define the matrix
\[
\Lambda = 
\begin{bmatrix}
\mu & 0 & 0 & \cdots & 0 \\
1 & \mu & 0 & \cdots & 0 \\
\vdots & \ddots & \ddots & \ddots & \vdots \\
0 & \cdots & 1 & \mu & 0 \\
0 & \cdots & 0 & 1 & \mu
\end{bmatrix}_{r\times r},
\]
so that
\[(f^*D_1,f^*D_2,\ldots,f^*D_r)=(D_1,D_2,\ldots,D_r)\Lambda.\]

Choose $L_i\in \mr{Pic}(X)_{\mb{R}}$ such that $p(L_i)=D_i$ for $i=1,2,\ldots,r$, and set $\mib{L}:=(L_1,L_2,\ldots,L_r)\in \mr{Pic}(X)_{\mb{R}}^{r}$. Let $V\subseteq\mr{Pic}(X)_{\mb{R}}$ be the finite dimensional $f^*$-invariant subspace generated by $L_1,\dots,L_r$. Then $\mib{M}:=f^*\mib{L}-\mib{L}\Lambda=(M_1,M_2,\ldots,M_r) \in (V\cap\mr{Pic}^0(X)_{\mb{R}})^r$.

We now show that the sequence $\big((f^*)^n\mib{L}\big)\Lambda^{-n}$ converges in $V^{r}$. Note that
\[
\begin{aligned}
&\big((f^*)^n\mib{L}\big)\Lambda^{-n}-\big((f^*)^{n-1}\mib{L}\big)\Lambda^{-(n-1)}\\
=&\big((f^{*})^{n-1}(f^*\mib{L}-\mib{L}\Lambda)\big)\Lambda^{-n}  \\
=&\big((f^{*})^{n-1}\mib{M}\big)\Lambda^{-n}. \\
\end{aligned}
\]

Notice that $\{(f^*)^n\mib{M}\}_{n\geq0}$ is a linear recurrence sequence, which admits the minimal polynomial of $f^*:\mr{Pic}^0(X)\rightarrow\mr{Pic}^0(X)$ as its characteristic polynomial. So the norm of $\big((f^{*})^{n-1}\mib{M}\big)\Lambda^{-n}$ is $O(n^{r+2g}(\frac{\rho_0}{\mu})^n)$, in which $g=\dim\Pic^0(X)$. It follows that the sequence $\big((f^*)^n\mib{L}\big)\Lambda^{-n}$ converges to some $\wt{\mib{L}}=(\wt{L}_1,\wt{L}_2,\ldots,\wt{L}_r)\in V^{r}$. By construction, we have $p(\wt{\mib{L}})=\mib{D}=(D_1,D_2,\ldots,D_r)\in\mr{N}^1(X)_{\mb{R}}^{r}$ and $f^*\wt{\mib{L}}=\wt{\mib{L}}\Lambda$.

Define $\iota(D_i):=\wt{L}_i, i=1,2,\ldots,r$. This extends linearly to a map $\iota: E\ra \mr{Pic}(X)_{\mb{R}}$. Repeating this construction for suitable cyclic subspaces of $\bigoplus\limits_{\mu>\rho_0}E_{\mu}$, we obtain the desired lifting $\iota:\bigoplus\limits_{\mu>\rho_0}E_{\mu} \ra \mr{Pic}(X)_{\mb{R}}$.
\end{proof}

\end{document}
